\documentclass[11pt]{letter}
\usepackage[margin=1in]{geometry}
\usepackage{hyperref}

\signature{Seongjin Choi, Ph.D.\\
Assistant Professor\\
Department of Civil, Environmental, and Geo-Engineering\\
University of Minnesota, Twin Cities\\
\href{mailto:choi@umn.edu}{choi@umn.edu}}

\address{Seongjin Choi\\
Department of Civil, Environmental, and Geo-Engineering\\
University of Minnesota, Twin Cities\\
500 Pillsbury Drive SE\\
Minneapolis, MN 55455, USA}

\date{\today}

\begin{document}

\begin{letter}{Editor-in-Chief\\
Transportation Research Part C: Emerging Technologies}

\opening{Dear Editor,}

We are pleased to submit our manuscript entitled \textbf{``Characterizing E-Scooter Riding Safety Through City-Scale Speed Profile Analysis''} for consideration as an Original Research Article in \textit{Transportation Research Part C: Emerging Technologies}.

E-scooter speeding is a growing safety concern worldwide, yet large-scale empirical evidence on speed behavior patterns remains scarce. This manuscript presents what we believe is the largest empirical analysis of e-scooter speed behavior to date, using 2.78 million GPS trajectories with per-point speed profiles from 382,328 unique users across 52 Korean cities.

Our study makes several contributions that we believe are well-suited for \textit{TR-C}:

\begin{enumerate}
    \item \textbf{Scale and data richness}: We analyze city-scale sensor-based speed profiles (not GPS-derived speeds), providing unprecedented granularity in characterizing e-scooter speed behavior.
    \item \textbf{Methodological breadth}: We combine spatial statistics (Getis-Ord $G_i^*$, Moran's $I$), logistic regression with GEE, mixed-effects models, a two-part hurdle model, and Gaussian mixture modeling to provide a comprehensive analysis.
    \item \textbf{Natural experiment}: We exploit a unique feature of the platform---a speed governor toggle---to estimate the causal effect of speed limiters through within-subject comparisons among 12,046 users, finding that the governor reduces speeding by 95\%.
    \item \textbf{Policy relevance}: Our findings provide direct empirical support for mandatory speed governors as the single most effective intervention, alongside geofencing and infrastructure-based strategies.
\end{enumerate}

The manuscript has not been published previously and is not under consideration elsewhere. All authors have approved the manuscript and agree with its submission to \textit{Transportation Research Part C}.

We believe this work advances the understanding of micromobility safety through rigorous, large-scale data analysis and offers actionable insights for regulators and operators. We look forward to your consideration.

\closing{Sincerely,}

\end{letter}
\end{document}
